\begin{myexercise}
	Erstellen Sie ein Wareninventarsystem, welches die Produkte deren Lagerbestand und Verkäufe in einem kleinen Online- oder physischen Geschäft verwaltet. 
	Erstellen Sie dafür folgende Klassen.
	\begin{itemize}
		\item Produktklasse - diese repräsentiert einen Artikel, der im Laden verkauft wird. \\Attribute: 
		\begin{itemize}
			\item name (string) 
			\item artikelnummer (string) 
			\item preis (float) 
			\item lagerbestand (integer)
		\end{itemize}
		Methoden:
		\begin{itemize}
			\item \_\_init\_\_(self, name, art\_nr, preis, bestand): Konstruktor.
			\item zeige\_details(self): Gibt Name, Preis und Bestand aus.
			\item bestand\_erhoehen(self, menge): Erhöht den Lagerbestand.
			\item bestand\_reduzieren(self, menge): Reduziert den Lagerbestand, muss prüfen, ob genug auf Lager ist.
		\end{itemize}
		\item Warenlager - diese Klasse verwaltet die Sammlung an Produkt-Objekten.\\Attribute:
		\begin{itemize}
			\item inventar (Dictionary, Schlüssel: artikelnummer, Wert: Produkt-Objekt).
			\item umsatz\_gesamt (float)
		\end{itemize}
		Methoden:
		\begin{itemize}
			\item \_\_init\_\_(self): Initialisiert inventar und umsatz\_gesamt.
			\item produkt\_hinzufuegen(self, produkt): Fügt ein Produkt-Objekt hinzu.
			\item produkt\_verkaufen(self, art\_nr, menge): Sucht das Produkt. Ruft die Methode bestand\_reduzieren() des Produkts auf. Wenn der Verkauf erfolgreich ist, erhöht es den umsatz\_gesamt der Warenlager-Klasse.
			\item zeige\_inventar(self): Listet alle Produkte und deren aktuellen Bestand auf.
			\item zeige\_umsatz(self): Gibt den aktuellen Gesamtumsatz aus.
		\end{itemize}
	\end{itemize}
	\begin{myanswer}
		\lstinputlisting[language=Python,caption=\texttt{produkt.py},label=listing:produkt,basicstyle=\ttfamily\scriptsize]{Praktikum/LPU/Code/produkt.py} 
	\end{myanswer}
	\begin{myanswer}
		\lstinputlisting[language=Python,caption=\texttt{warenlager.py},label=listing:warenlager,basicstyle=\ttfamily\scriptsize]{Praktikum/LPU/Code/warenlager.py}
	\end{myanswer}
\end{myexercise}